% Regression: t2_pentagon2 — renders the MPO-wire fusion move c -> (a, b).
% Target: RMP 2011.12127 fig3_pentagon2.pdf (MPO fat-graph fusion vertex).
% Formula: the MPO-wire fusion move  c -> (a, b):
%   red directed MPO wires  c: n -> m,  a: m -> j,  b: m -> k,
%   fat-strip faces A (above a), B (the cul-de-sac between a and b),
%   C (below b); gray wedges at the wire endpoints n, m, j, k.
% Genuinely non-grid (trivalent junction of virtual wires): tenkzfree.
% Honest gaps: no fat-strip/ribbon face outline primitive and no
% free-tier annotation label, so the faces A, B, C are absent; no
% gray-wedge terminal skin (house beads instead); \tnjoin has no wire-hue
% key, so the red MPO hue renders as house ink.
\documentclass[varwidth=19cm, border=6pt]{standalone}
\usepackage{amsmath, amssymb}
\usepackage{tenkz}
\begin{document}

\[
  \begin{tenkzfree}
    % wire endpoints (paper: gray wedges), typed virtual on the used anchors
    \tnput[dot, label pos=west, ports={east:virtual}]{n}{(0,0)}{n}
    \tnput[dot, label pos=south,
           ports={west:virtual, north:virtual, south:virtual}]{m}{(16mm,0)}{m}
    \tnput[dot, label pos=east, ports={west:virtual}]{j}{(30mm,9mm)}{j}
    \tnput[dot, label pos=east, ports={west:virtual}]{k}{(30mm,-9mm)}{k}
    % c: n -> m (straight), then the split a: m -> j and b: m -> k (curved)
    \tnjoin[dir, label=c]{n.east}{m.west}
    \tnjoin[dir, route=curve, out=50, in=195, label=a]{m.north}{j.west}
    \tnjoin[dir, route=curve, out=-50, in=165, label=b]{m.south}{k.west}
  \end{tenkzfree}
\]

\end{document}
